\chapter*{Conclusion and Perspectives}
\addcontentsline{toc}{chapter}{Conclusion and Perspectives}

This project was carried out as part of my final-year internship at FinGenesis, addressing the operational challenges posed by fragmented infrastructure management tools. Engineers previously relied on multiple disconnected platforms—Grafana for monitoring, GitLab for repository management, AWS Console for cloud resources, and manual communication channels. This fragmentation resulted in time-consuming workflows, limited visibility, and scalability issues.

The Infrastructure Command Center was developed as a centralized platform consolidating these operations into a single, user-friendly interface. The platform delivers three core features: Dynamic Visual Workflows for infrastructure visualization, Control Panel for repository and deployment management, and ML Metrics Monitoring for model performance tracking.

The backend was built using Python and FastAPI with asynchronous operations and intelligent caching, achieving dramatic performance improvements—reducing load times from 30+ seconds to under 50 milliseconds. The frontend uses React, ReactFlow, and Recharts, delivering a responsive interface. The application is containerized with Docker and deployed on AWS using Kubernetes, with automated CI/CD through GitLab.

Two major technical challenges were successfully resolved: implementing React Flow for maintainable workflow visualization, and optimizing GitLab API performance through async/await patterns and caching, achieving 99\% performance improvement.

As a junior developer, this project provided intensive learning across the full technology stack—React, FastAPI, Docker, Kubernetes, and GitLab CI/CD—while working in a professional environment with enterprise-level standards. The Infrastructure Command Center achieved its objective of transforming fragmented infrastructure management into a streamlined operation. FinGenesis engineers have begun adopting the platform in their workflows, and it will continue evolving based on operational needs.

\section*{Perspectives}

Several enhancements can extend the platform's capabilities:

\textbf{Scheduled Job Management:} Add visibility and control over scheduled jobs, cron tasks, and automated processes for complete infrastructure oversight.

\textbf{Enhanced Notifications:} Implement real-time alerts for pipeline failures, deployment completions, and ML model drift, integrated with Discord or Slack.

\textbf{Business Development Features:}
\begin{itemize}
    \item \textbf{Expense Tracking}: Display AWS costs, API subscriptions, and infrastructure spending for cost optimization
    \item \textbf{Data Policy Management}: Centralize documentation of data sources, backups, and retention policies
\end{itemize}

\textbf{Knowledge Base:}
\begin{itemize}
    \item \textbf{Technology Stack Visualization}: Display used technologies, AWS services, and libraries in an interactive format
    \item \textbf{Package Documentation}: Integrate auto-generated documentation for internal packages
\end{itemize}

\textbf{Additional Enhancements:}
\begin{itemize}
    \item Role-based access control for sensitive operations
    \item Audit logging for compliance and security
    \item Mobile optimization for on-the-go monitoring
    \item Extended monitoring system integration
\end{itemize}

These enhancements will make the Infrastructure Command Center more comprehensive and scalable, aligning with FinGenesis's long-term operational goals. The platform's modular architecture supports these additions without requiring fundamental redesigns.